\subsection{Programación de Intervalos con Cobertura por Grupos (DP de Bitmask + Fenwick de Máximos)}

\graybold{Preguntas guía:}

\begin{itemize}
\item ¿Hay que elegir un subconjunto de intervalos mutuamente compatibles (no solapados) que maximice un valor total?
\item ¿Además existe una restricción de cobertura: cada uno de varios grupos debe tener AL MENOS un intervalo elegido?
\item ¿El número de grupos es pequeño (≤ \textasciitilde{}20, para caber en una máscara de bits), aunque el número total de intervalos sea grande?
\end{itemize}

\graybold{¿Para qué sirve?} Generalizar la programación de intervalos ponderada (weighted interval scheduling) clásica al caso donde los intervalos están agrupados y se exige cubrir TODOS los grupos con al menos un intervalo, maximizando el valor total de los elegidos (pudiendo elegir más de uno por grupo si conviene).

\graybold{Complejidad:} \bigO{M \cdot 2^G \cdot log M}, con M = número total de intervalos y G = número de grupos.

\graybold{Fundamento teórico:} Se procesan los intervalos ordenados por tiempo de FIN (como en programación de intervalos clásica). Para cada máscara de grupos "ya cubiertos", se mantiene un Fenwick de MÁXIMOS indexado por tiempo de fin (con coordenadas comprimidas), que responde "¿cuál es el mejor valor acumulado usando intervalos compatibles con un tiempo de inicio dado?" en \bigO{log M}. Al procesar cada intervalo, se combina con el mejor resultado compatible de la máscara SIN su grupo (o la misma máscara, si el grupo ya estaba cubierto), y se actualiza la máscara resultante — igual que un DP de subconjuntos, pero usando un Fenwick en vez de un arreglo simple para manejar la dimensión de compatibilidad temporal.

\graybold{Ejercicio aplicado}

Un festival tiene N escenarios (grupos), cada uno con varios shows (intervalos de tiempo con un valor asociado). Hay que ver AL MENOS un show en cada escenario, y los shows elegidos (de cualquier escenario) deben ser mutuamente compatibles en el tiempo. Maximizar el valor total, o reportar -1 si no hay forma de cubrir todos los escenarios.

\graybold{I/O:} N ≤ 10 (hasta 1024 máscaras), suma de shows ≤ 1000 → lectura total con tokenización (patrón 2); el Fenwick de máximos se implementa con arreglos planos (no clases ni closures) para evitar el overhead de Python en el doble loop show×máscara.

\graybold{Código solución}

\begin{lstlisting}[language=Python]
import sys
from bisect import bisect_right

NEG = -1

def solve():
    data = sys.stdin.buffer.read().split()
    idx = 0
    n = int(data[idx]); idx += 1
    shows = []
    for s in range(n):
        mi = int(data[idx]); idx += 1
        for _ in range(mi):
            a = int(data[idx]); b = int(data[idx + 1]); c = int(data[idx + 2]); idx += 3
            shows.append((b, a, c, s))   # ordenar por fin: b primero

    shows.sort()

    end_times = sorted(set(x[0] for x in shows))
    sz = len(end_times)
    rank = {t: i + 1 for i, t in enumerate(end_times)}

    num_masks = 1 << n
    full_mask = num_masks - 1
    # un Fenwick de maximos por mascara, todo en una lista plana (mask*(sz+1)+pos)
    fen = [NEG] * (num_masks * (sz + 1))

    for (b, a, c, s) in shows:
        pos = bisect_right(end_times, a)   # posicion del mayor fin <= inicio de este show
        bit = 1 << s
        rb = rank[b]
        for mask in range(num_masks):
            if mask == 0:
                best_prev = 0              # "no elegir nada aun" siempre es compatible
            elif pos == 0:
                continue
            else:
                base = mask * (sz + 1)
                res = NEG
                i = pos
                while i > 0:                # query de maximo de prefijo, inline
                    v = fen[base + i]
                    if v > res:
                        res = v
                    i &= i - 1
                if res == NEG:
                    continue                # esta mascara nunca fue alcanzable aun
                best_prev = res
            new_mask = mask | bit           # si el grupo ya estaba, new_mask == mask
            val = best_prev + c
            base2 = new_mask * (sz + 1)
            i = rb
            while i <= sz:                  # update de maximo, inline
                if val > fen[base2 + i]:
                    fen[base2 + i] = val
                i += i & (-i)

    base = full_mask * (sz + 1)
    res = NEG
    i = sz
    while i > 0:
        v = fen[base + i]
        if v > res:
            res = v
        i &= i - 1
    print(res if res != NEG else -1)

solve()
\end{lstlisting}
